\documentclass[9pt,a4paper,twocolumn,twoside]{tau-class/tau}
\usepackage[spanish,es-nodecimaldot,es-noindentfirst]{babel}

\doctype{Reporte Ejecutivo Integrado}
\title{Reporte de performance y panorama competitivo: Centro de Psicologia Integral y Terapia - Psicotelcon}

\author[a,1]{Equipo de Analitica Oderbiz}
\affil[a]{Unidad de Inteligencia Comercial}

\dates{Compilado el \today}
\docinfo{Documento generado con datos internos en formato \texttt{llm\_context\_report.page.v1} e investigacion competitiva publica. Toda lectura no directamente observable se etiqueta como \texttt{inferencia}.}
\footinfo{Reporte Oderbiz}
\theday{\today}
\organization{Oderbiz}
\leadauthor{Equipo Oderbiz}

\begin{abstract}
En la ventana \texttt{last\_30d}, la cuenta analizada invirtio \$21.30 y alcanzo 14{,}759 impresiones con 12{,}695 de alcance (\texttt{observado}). El rendimiento de atencion inicial es positivo (CTR 1.57\% y 209 clics unicos), con conversion conversacional efectiva de 11 conversaciones iniciadas y 11 primeras respuestas (\texttt{observado}). Sin embargo, la calidad de trafico muestra dependencia total de mensajeria dentro de plataforma al no registrar outbound clicks (\texttt{observado}). El entorno competitivo local presenta alta fragmentacion entre profesionales individuales, centros integrales y directorios comparativos (\texttt{observado}). La oportunidad prioritaria es fortalecer diferenciacion por problema clinico, estandarizar arquitectura de campanas y medir calidad post-conversacion para mejorar eficiencia de crecimiento (\texttt{inferencia}).
\end{abstract}

\keywords{Meta Ads, psicologia, Loja, mensajeria, inteligencia competitiva}

\begin{document}
\maketitle
\thispagestyle{firststyle}

\section{Resumen ejecutivo}

La operacion paid actual genera volumen eficiente en alcance y clics con baja inversion absoluta (\texttt{observado}). El cuello de botella principal no es la captacion de atencion sino la trazabilidad de negocio posterior al primer contacto (\texttt{inferencia}). En geo, Loja concentra casi todo el resultado (10 de 11 conexiones de mensajeria) y Zamora Chinchipe actua como extension secundaria (\texttt{observado}). En demografia, 25--44 concentra volumen de conversion conversacional y 55+ aporta picos de CTR/costo por clic favorables en menor escala (\texttt{observado}+\texttt{inferencia}). Se recomienda ejecutar un plan escalonado P1/P2/P3 con foco en estructura de cuenta, creatividad por dolor y medicion de calidad real de lead.

\section{Contexto del negocio y periodo analizado}

\begin{itemize}
  \item Nombre de pagina: \texttt{Centro de Psicologia Integral y Terapia - Psicotelcon}
  \item ID de pagina: \texttt{1506380769434870}
  \item Nombre de cuenta: \texttt{Dra Marizol Jimenez Luzon}
  \item ID de cuenta: \texttt{act\_131112367482947}
  \item Campana filtrada: \texttt{N/A} (sin filtro aplicado)
  \item ID campana filtrada: \texttt{N/A}
  \item Moneda: \texttt{USD}
  \item Zona horaria: \texttt{America/Guayaquil}
  \item Rango analizado: \texttt{last\_30d} (date\_range sin \texttt{since/until} explicitos)
\end{itemize}

La cuenta mantiene un inventario historico amplio de campanas activas con objetivos dominantes de \texttt{MESSAGES} y \texttt{OUTCOME\_ENGAGEMENT}, con presencia puntual de \texttt{OUTCOME\_TRAFFIC}/\texttt{LINK\_CLICKS} (\texttt{observado}).

\section{Baseline interno}

\subsection{KPIs base}
\begin{itemize}
  \item Inversion: \$21.30; impresiones: 14{,}759; alcance: 12{,}695 (\texttt{observado}).
  \item CTR: 1.565\%; CPM: \$1.44 (\texttt{observado}).
  \item Embudo: 209 clics unicos; 11 conversaciones iniciadas; 11 primeras respuestas (\texttt{observado}).
  \item Calidad de trafico: 0 outbound clicks; costo por clic unico: \$0.10 (\texttt{observado}).
\end{itemize}

\subsection{Tendencia temporal}
Los dias con mayor entrega y gasto (25-mar, 20-abr, 21-abr) sostienen CTR entre 1.56\% y 1.88\%, con variacion de CPM en rango 1.19--1.67 (\texttt{observado}). Esto sugiere un sistema estable a pequena escala, pero con sensibilidad a subasta cuando aumenta la presion de entrega diaria (\texttt{inferencia}).

\subsection{Senal de embudo y calidad}
La relacion entre 209 clics unicos y 11 conversaciones iniciadas implica una tasa de paso aproximada de 5.3\% de clic unico a conversacion (\texttt{estimado}). La ausencia de outbound clicks confirma que el embudo actual se resuelve principalmente en entorno Meta, reduciendo visibilidad de conversion fuera de plataforma (\texttt{observado}+\texttt{inferencia}).

\subsection{Senal geo y demografica}
\textbf{Geo}: Loja concentra 96.4\% del gasto (20.53/21.30) y 10/11 resultados de mensajeria, mientras Zamora Chinchipe aporta escala marginal con costo por resultado menor (\$0.77 vs \$2.05) sobre muestra pequena (\texttt{observado}+\texttt{inferencia}).

\textbf{Edad}: los rangos 25--44 concentran volumen de impresiones y resultados de mensajeria; 55--64 y 65+ muestran CTR alto y CPC bajo en menor escala, potencialmente utiles para pruebas de eficiencia incremental (\texttt{observado}+\texttt{inferencia}).

\section{Panorama competitivo}

\subsection{Lectura del entorno}
El mercado local combina:
\begin{itemize}
  \item profesionales individuales con fuerte posicionamiento por especialidad;
  \item centros con propuesta integral (clinica, educativa, lenguaje, familia);
  \item directorios/plataformas que facilitan comparacion por perfil, modalidad y reputacion.
\end{itemize}

Este contexto incrementa la necesidad de diferenciacion explicita por resultados terapeuticos, no solo por disponibilidad de cita (\texttt{inferencia}).

\subsection{Fuentes competitivas publicas consultadas}
\begin{itemize}
  \item \url{https://psicologiaymente.com/psicologos/2069561/dra-marizol-jimenez-luzon}
  \item \url{https://psicologiaymente.com/psicologos/2062410/nicole-loaiza}
  \item \url{https://www.psico.org/ec/directorio/loja}
  \item \url{https://www.psico.org/ec/loja/terapia-de-pareja}
  \item \url{https://www.terappio.com/ec/family-therapist/ecuador/loja}
  \item \url{https://tecnologicosudamericano.edu.ec/bienestar-estudiantil/centro-psicoterapeutico-la-cabana/}
\end{itemize}

\section{Diagnostico integrado}

\subsection{Fortalezas}
\begin{itemize}
  \item Eficiencia de cobertura a bajo presupuesto, con CPM competitivo (\texttt{observado}).
  \item Capacidad de activar conversacion directa (11/11 inicio-respuesta) (\texttt{observado}).
  \item Presencia continua de campanas y biblioteca amplia de activos (\texttt{observado}).
\end{itemize}

\subsection{Brechas}
\begin{itemize}
  \item Cero salida outbound y baja trazabilidad de conversion final (\texttt{observado}).
  \item Nomenclatura heterogenea de campanas que dificulta aprendizaje acumulado (\texttt{observado}).
  \item Cobertura de diagnostico por anuncio aun limitada a filas agregadas (\texttt{observado}+\texttt{inferencia}).
\end{itemize}

\subsection{Riesgos}
\begin{itemize}
  \item Dependencia de un canal de conversacion sin capa CRM cerrada (\texttt{observado}+\texttt{inferencia}).
  \item Vulnerabilidad a comparacion en directorios y oferta local fragmentada (\texttt{observado}).
  \item Variabilidad de costo al escalar presupuesto diario (\texttt{inferencia}).
\end{itemize}

\subsection{Oportunidades}
\begin{itemize}
  \item Creatividades segmentadas por motivo de consulta (lenguaje, conducta, desarrollo, familia) (\texttt{inferencia}).
  \item Embudo dual \texttt{MESSAGES + TRAFFIC} para mejorar atribucion (\texttt{inferencia}).
  \item Microexpansion geografica controlada en provincias adyacentes con CPA competitivo (\texttt{inferencia}).
\end{itemize}

\section{Plan de accion priorizado (P1/P2/P3)}

\subsection{P1: alto impacto / bajo esfuerzo}
\begin{itemize}
  \item Estandarizar naming: objetivo--audiencia--oferta--fecha.
  \item Definir tablero semanal minimo: costo por conversacion iniciada, costo por primera respuesta y costo por cita agendada.
  \item Crear 3 variantes de copy por dolor especifico con CTA a WhatsApp.
  \item Establecer protocolo operativo de respuesta para reducir tiempo a primer contacto.
\end{itemize}

\subsection{P2: alto impacto / esfuerzo medio-alto}
\begin{itemize}
  \item Activar campana complementaria de trafico con landing de conversion.
  \item Implementar taxonomia de audiencias por edad del hijo y necesidad terapeutica.
  \item Diseñar test de creatividades por formato (testimonio, educativo, caso-clinico).
  \item Integrar medicion post-chat (contacto efectivo, asistencia, continuidad).
\end{itemize}

\subsection{P3: experimental}
\begin{itemize}
  \item Lookalikes de conversaciones calificadas.
  \item Experimentacion con contenido de autoridad local y alianzas.
  \item Secuencias de retargeting segun profundidad conversacional.
\end{itemize}

\section{Fuentes y limitaciones}

\subsection{Fuentes internas}
\begin{itemize}
  \item \texttt{report\_metadata}
  \item \texttt{page\_overview}
  \item \texttt{timeseries\_daily}
  \item \texttt{funnel}
  \item \texttt{traffic\_quality}
  \item \texttt{geo}
  \item \texttt{demographics}
  \item \texttt{campaigns\_available}
  \item \texttt{reference\_link\_coverage}
  \item \texttt{module\_status}
\end{itemize}

\subsection{Fuentes visuales}
No se incorporaron capturas en esta corrida automatizada. La evidencia competitiva se documento con URLs trazables y fecha de generacion del reporte (\texttt{2026-04-25}) como fallback parcial.

Rutas esperadas para proxima iteracion visual:
\begin{itemize}
  \item \texttt{docs/reports/figures/centro-de-psicologia-integral-y-terapia-psicotelcon/}
  \item Convencion: \texttt{fig-<seccion>-<fuente>-<YYYYMMDD>-vN.png}
\end{itemize}

\subsection{Limitaciones detectadas}
\begin{itemize}
  \item Sin token/credito confirmado para \texttt{apify-market-research} en esta ejecucion.
  \item Sin captura automatica de evidencias web/redes en esta corrida; se aplico fallback textual con fuentes URL.
  \item Sin datos CRM de cierre (citas atendidas, ticket, recompra), por lo que la lectura de impacto en ingreso es parcial.
\end{itemize}

\end{document}
